% ========================================
% PREAMBLE for Probability and Statistics
% ========================================

% --- Basics ---
\usepackage[utf8]{inputenc}
\usepackage[T1]{fontenc}
\usepackage{textcomp}
\usepackage{graphicx}
\usepackage{float}
\usepackage{booktabs}
\usepackage{enumitem}
\usepackage{emptypage}
\usepackage{subcaption}
\usepackage{multicol}
\usepackage[usenames,dvipsnames]{xcolor}
\setlength{\headheight}{14.49998pt}

% --- Math Packages ---
\usepackage{amsfonts, mathtools, amsthm, amssymb} %% REMOVED: amsmath (loaded by mathtools)
\usepackage{mathrsfs}
\usepackage{cancel}
\usepackage{bm}
\usepackage[cal=boondoxo]{mathalfa}
\usepackage{systeme}
\usepackage{stmaryrd}
\usepackage{enumitem}

% --- Math Commands ---
\newcommand\N{\ensuremath{\mathbb{N}}}
\newcommand\R{\ensuremath{\mathbb{R}}}
\newcommand\Z{\ensuremath{\mathbb{Z}}}
\renewcommand\O{\ensuremath{\emptyset}}
\newcommand\Q{\ensuremath{\mathbb{Q}}}
\newcommand\C{\ensuremath{\mathbb{C}}}
\newcommand{\samplespace}{\mathcal{S}}
\DeclareMathOperator{\sgn}{sgn}
%% REMOVED: Global redefinition of \lim to use \limits. This is poor practice as it disrupts inline text flow. Use \lim\limits or \displaystyle where needed.
\let\implies\Rightarrow
\let\impliedby\Leftarrow
\let\iff\Leftrightarrow
\let\epsilon\varepsilon
\newcommand\contra{\scalebox{1.1}{$\lightning$}}

% --- Correction Commands ---
\definecolor{correct}{HTML}{009900}
\newcommand\correct[2]{\ensuremath{\:}{\color{red}{#1}}\ensuremath{\to }{\color{correct}{#2}}\ensuremath{\:}}
\newcommand\green[1]{{\color{correct}{#1}}}

% --- Utilities ---
\usepackage{xifthen} %% MOVED: xifthen here to ensure it's loaded before being used.
\newcommand\hr{\noindent\rule[0.5ex]{\linewidth}{0.5pt}}
\newcommand\hide[1]{}
\newcommand\mat[1]{\mathbf{#1}}

% --- SI Units ---
\usepackage{siunitx}
\sisetup{locale = FR}

% --- TikZ ---
\usepackage{tikz}
\usepackage{tikz-cd}
\usetikzlibrary{intersections, angles, quotes, calc, positioning}
\usetikzlibrary{arrows.meta}
\usetikzlibrary{shapes,positioning,trees,calc}
\usepackage{pgfplots}
\pgfplotsset{compat=1.13}

\tikzset{
    force/.style={thick, {Circle[length=2pt]}-Stealth, shorten <=-1pt} %% FIXED: 'stealth' should be 'Stealth' (case-sensitive)
}

% --- Page Layout ---
\usepackage[
    top=3cm,
    bottom=3cm,
    inner=3cm,
    outer=2.5cm
]{geometry}
\setlength{\headheight}{13.6pt} %% MOVED: Set headheight after loading geometry and fancyhdr for proper calculation.

% --- Theorem Environments ---
\usepackage{thmtools}
\usepackage[framemethod=TikZ]{mdframed}
\mdfsetup{skipabove=1em,skipbelow=0em}

\theoremstyle{definition}

\declaretheoremstyle[
    headfont=\bfseries\sffamily\color{ForestGreen!70!black}, bodyfont=\normalfont,
    mdframed={
        linewidth=2pt,
        rightline=false, topline=false, bottomline=false,
        linecolor=ForestGreen, backgroundcolor=ForestGreen!5,
    }
]{thmgreenbox}

\declaretheoremstyle[
    headfont=\bfseries\sffamily\color{NavyBlue!70!black}, bodyfont=\normalfont,
    mdframed={
        linewidth=2pt,
        rightline=false, topline=false, bottomline=false,
        linecolor=NavyBlue, backgroundcolor=NavyBlue!5,
    }
]{thmbluebox}

\declaretheoremstyle[
    headfont=\bfseries\sffamily\color{NavyBlue!70!black}, bodyfont=\normalfont,
    mdframed={
        linewidth=2pt,
        rightline=false, topline=false, bottomline=false,
        linecolor=NavyBlue
    }
]{thmblueline}

\declaretheoremstyle[
    headfont=\bfseries\sffamily\color{RawSienna!70!black}, bodyfont=\normalfont,
    mdframed={
        linewidth=2pt,
        rightline=false, topline=false, bottomline=false,
        linecolor=RawSienna, backgroundcolor=RawSienna!5,
    }
]{thmredbox}

\declaretheoremstyle[
    headfont=\bfseries\sffamily\color{RawSienna!70!black}, bodyfont=\normalfont,
    numbered=no,
    mdframed={
        linewidth=2pt,
        rightline=false, topline=false, bottomline=false,
        linecolor=RawSienna, backgroundcolor=RawSienna!1,
    },
    qed=\qedsymbol
]{thmproofbox}

\declaretheoremstyle[
    headfont=\bfseries\sffamily\color{NavyBlue!70!black}, bodyfont=\normalfont,
    numbered=no,
    mdframed={
        linewidth=2pt,
        rightline=false, topline=false, bottomline=false,
        linecolor=NavyBlue, backgroundcolor=NavyBlue!1,
    },
]{thmexplanationbox}


% --- Solution Environment ---
\declaretheoremstyle[
    headfont=\bfseries\sffamily\color{ForestGreen!70!black},
    bodyfont=\normalfont,
    numbered=no,
    mdframed={
        linewidth=2pt,
        rightline=false, topline=false, bottomline=false,
        linecolor=ForestGreen, backgroundcolor=ForestGreen!1,
    },
    qed=\qedsymbol
]{thmsolutionbox}

\declaretheorem[style=thmsolutionbox, name=Solution]{solution}


\declaretheorem[style=thmgreenbox, name=Definition]{definition}
\declaretheorem[style=thmbluebox, numbered=no, name=Example]{eg}
\declaretheorem[style=thmredbox, name=Proposition]{prop}
\declaretheorem[style=thmredbox, name=Theorem]{theorem}
\declaretheorem[style=thmredbox, name=Lemma]{lemma}
\declaretheorem[style=thmredbox, numbered=no, name=Corollary]{corollary}

%% ADDED: New 'exercise' environment using the 'thmbluebox' style. It will be numbered automatically.
\declaretheorem[style=thmbluebox, name=Exercise]{exercise}

\declaretheorem[style=thmproofbox, name=Proof]{replacementproof}
\renewenvironment{proof}[1][\proofname]{\vspace{-10pt}\begin{replacementproof}}{\end{replacementproof}}

\declaretheorem[style=thmexplanationbox, name=Proof]{tmpexplanation}
\newenvironment{explanation}[1][]{\vspace{-10pt}\begin{tmpexplanation}}{\end{tmpexplanation}}

\declaretheorem[style=thmblueline, numbered=no, name=Remark]{remark}
\declaretheorem[style=thmblueline, numbered=no, name=Note]{note}

\newtheorem*{notation}{Notation}
\newtheorem*{previouslyseen}{As previously seen}
\newtheorem*{problem}{Problem}
\newtheorem*{observe}{Observe}
\newtheorem*{property}{Property}
\newtheorem*{intuition}{Intuition}


% --- Lecture & Exercise Commands ---
\makeatletter
%% REMOVED: Old \exercise and \subexercise commands are now replaced by the 'exercise' environment below.
%% The new environment provides consistent styling and automatic numbering.

%% ADDED: A custom list for sub-exercises (parts a, b, c, ...). Use it inside the new 'exercise' environment.
\newlist{subtasks}{enumerate}{1}
\setlist[subtasks,1]{label=\alph*)}

\def\@lesson{}%
\newcommand{\lesson}[3]{
    \ifthenelse{\isempty{#3}}{%
        \def\@lesson{Lecture #1}%
    }{%
        \def\@lesson{Lecture #1: #3}%
    }%
    \subsection*{\@lesson}
    \marginpar{\small\textsf{\mbox{#2}}}
}
\makeatother


% --- Headers and Footers ---
\usepackage{fancyhdr}
\pagestyle{fancy}
\fancyhead{} %% Clears existing header/footer settings
\fancyhead[RO,LE]{\@lesson}
\fancyfoot[LE,RO]{\thepage}
\fancyfoot[C]{\leftmark}
\renewcommand{\headrulewidth}{0.4pt}


% --- Notes ---
\usepackage{todonotes}
\usepackage{tcolorbox}
\tcbuselibrary{breakable}

\newenvironment{verbetering}{\begin{tcolorbox}[
    arc=0mm,
    colback=white,
    colframe=green!60!black,
    title=Remark,
    fonttitle=\sffamily,
    breakable
]}{\end{tcolorbox}}

\newenvironment{noot}[1]{\begin{tcolorbox}[
    arc=0mm,
    colback=white,
    colframe=white!60!black,
    title=#1,
    fonttitle=\sffamily,
    breakable
]}{\end{tcolorbox}}



% --- Figure Support ---
\usepackage{import}
\pdfminorversion=7
\usepackage{pdfpages}
\usepackage{transparent}
\newcommand{\incfig}[1]{%
    \def\svgwidth{\columnwidth}
    \import{./figures/}{#1.pdf_tex}
}
\pdfsuppresswarningpagegroup=1

% --- Hyperlinks ---
%% MOVED: hyperref is best loaded late in the preamble.
%% REMOVED: url package (functionality is included in hyperref).
\usepackage{hyperref}
\hypersetup{
    colorlinks=false,
    pdfborder={0 0 0}
}